\subsection{Algorithm Specification}
\label{sec:algorithm-spec}

This section provides complete specification of the recursive bisection algorithm used to generate algorithmic redistricting plans, enabling replication and validation of our baseline efficiency gap estimates.

\subsubsection{Core Algorithm: Recursive Bisection via METIS}

We employ recursive bisection~\cite{karypis1998metis,deluca2026recursive}, a divide-and-conquer approach that hierarchically partitions census tracts into districts. The algorithm recursively splits geographic regions into two balanced subregions until each subregion contains exactly one district's worth of population.

\textbf{Algorithm pseudocode:}
\begin{verbatim}
recursive_bisection(tracts, adjacency_graph, target_districts):
    if target_districts == 1:
        return all tracts assigned to district 1

    # Split into two groups
    split_point = target_districts // 2
    partition = metis_bisect(adjacency_graph, tract_populations)

    # Separate into two subgraphs
    group_0 = tracts where partition[i] == 0
    group_1 = tracts where partition[i] == 1

    # Recursively partition each group
    assignments_0 = recursive_bisection(group_0, split_point)
    assignments_1 = recursive_bisection(group_1,
                                      target_districts - split_point)

    # Offset group_1 district IDs and combine
    return assignments_0 + assignments_1 (with offset)
\end{verbatim}

\textbf{Rationale for recursive bisection:} Binary splits are computationally efficient (requiring $\log_2 K$ recursion levels for $K$ districts) and produce well-balanced partitions at each level. Direct $K$-way partitioning struggles with large $K$ values and exhibits poor population balance~\cite{karypis1998metis}.

\subsubsection{METIS Parameterization}

At each recursion level, we invoke METIS (Multi-constraint Graph Partitioning library) to perform the binary partition. METIS configuration critically affects compactness and partisan outcomes.

\textbf{Edge weight specification:}
\begin{itemize}
    \item \textbf{Vertex weights}: Census tract populations (to enforce equal-population districts)
    \item \textbf{Edge weights}: Uniform (unweighted edges between adjacent tracts)
    \item \textbf{Justification}: Geographic compactness emerges from minimizing edge cuts without introducing additional geographic biases through distance weighting
\end{itemize}

\textbf{METIS parameters:}
\begin{itemize}
    \item \texttt{nparts = 2}: Binary partition at each recursion level
    \item \texttt{niter = 100}: Number of refinement iterations (higher = better quality, slower)
    \item \texttt{ufactor = 10}: Population imbalance tolerance ($\pm$0.5\% of target)
    \item \texttt{objtype = 'cut'}: Minimize edge cut (default METIS objective)
    \item \texttt{seed = deterministic}: Fixed random seed for reproducibility
\end{itemize}

\textbf{Tiebreaking:} When multiple partitions achieve similar edge cuts, METIS uses its internal refinement heuristics (Kernighan-Lin algorithm) to select among near-optimal solutions. We do not intervene in this process, maintaining algorithmic neutrality.

\subsubsection{Population Balance Enforcement}

Constitutional requirements mandate districts within $\pm$0.5\% of ideal population~\cite{karcher1993right}. We enforce this through METIS's vertex weighting mechanism:

\begin{enumerate}
    \item \textbf{Target population}: Total state population $\div$ number of districts
    \item \textbf{Tolerance}: $\pm$0.5\% of target (ufactor = 10 in METIS)
    \item \textbf{Hierarchical balancing}: At each recursion level, ensure subregions are balanced before further splitting
    \item \textbf{Verification}: Post-processing check confirms all districts meet $\pm$0.5\% threshold
\end{enumerate}

\textbf{Balancing result:} All 435 algorithmic districts across 50 states achieve population deviations $<$0.3\% of target (mean absolute deviation: 0.12\%), well within constitutional requirements.

\subsubsection{Partisan Data Exclusion}

Critical to our neutrality claim: the algorithm cannot access partisan data during redistricting.

\textbf{Input data:}
\begin{itemize}
    \item Census tract boundaries (geographic polygons)
    \item Census tract populations (2020 Census)
    \item Tract adjacency relationships (computed from spatial boundaries)
\end{itemize}

\textbf{Explicitly excluded:}
\begin{itemize}
    \item Election results (precinct-level or block-level)
    \item Voter registration data
    \item Partisan indices or proxies
    \item Demographic variables correlated with partisanship (race, income, education)
\end{itemize}

\textbf{Verification:} We confirm partisan data exclusion by regenerating algorithmic plans with identical METIS parameters but different random seeds. If algorithms accessed partisan data, seed changes would not affect outcomes; observed variation in district boundaries (with stable efficiency gaps) confirms genuine algorithmic neutrality.

\subsubsection{Sensitivity Analysis}

To validate that our -3.2\% efficiency gap baseline is not METIS-specific, we test sensitivity to key algorithmic choices for five representative states (Pennsylvania, Texas, California, Florida, North Carolina):

\textbf{Alternative algorithms tested:}
\begin{itemize}
    \item \textbf{K-means clustering}: Districts formed by assigning tracts to nearest district centroid (k=435 for national analysis)
    \item \textbf{Shortest splitline}: Recursive splitting along shortest geographic lines bisecting population
    \item \textbf{Voronoi diagrams}: Districts as Voronoi cells around population-weighted seed points
\end{itemize}

\textbf{Edge weight variations:}
\begin{itemize}
    \item \textbf{Unweighted} (baseline): Equal edge weights
    \item \textbf{Inverse distance}: $w_{ij} = 1 / d_{ij}$ where $d_{ij}$ is centroid distance
    \item \textbf{Shared boundary length}: $w_{ij} = $ length of shared boundary between tracts $i$ and $j$
\end{itemize}

\textbf{Sensitivity results:}
\begin{itemize}
    \item Alternative algorithms produce efficiency gaps ranging from -2.8\% to -3.6\% (mean: -3.1\%, std dev: 0.3\%)
    \item Edge weight variations produce efficiency gaps from -3.0\% to -3.4\% (mean: -3.2\%, std dev: 0.2\%)
    \item All variations remain substantially below enacted plan efficiency gaps (+5.1\%)
\end{itemize}

\textbf{Interpretation:} The -3.2\% baseline is robust across neutral algorithmic approaches. Minor variations reflect different compactness definitions, but all neutral algorithms cluster within a narrow band ($\pm$0.4 percentage points), far from enacted plan bias.

\subsubsection{Ensemble Generation and Uncertainty Quantification}

For the five sensitivity analysis states, we generated ensembles of 100 algorithmic plans per state by varying the random seed while holding all other METIS parameters constant. This quantifies inherent uncertainty in algorithmic redistricting.

\textbf{Ensemble statistics (5 states, 100 maps each):}
\begin{itemize}
    \item \textbf{Pennsylvania}: Mean EG = -2.8\%, std dev = 0.3\%, range = [-3.4\%, -2.3\%]
    \item \textbf{Texas}: Mean EG = -3.6\%, std dev = 0.4\%, range = [-4.3\%, -2.8\%]
    \item \textbf{California}: Mean EG = -3.1\%, std dev = 0.2\%, range = [-3.6\%, -2.7\%]
    \item \textbf{Florida}: Mean EG = -3.2\%, std dev = 0.3\%, range = [-3.9\%, -2.6\%]
    \item \textbf{North Carolina}: Mean EG = -3.0\%, std dev = 0.3\%, range = [-3.7\%, -2.4\%]
\end{itemize}

\textbf{Key finding:} Within-state variation (std dev $\approx$ 0.3\%) is substantially smaller than the algorithmic-enacted gap (8.3 percentage points). This confirms our single algorithmic plan per state provides a stable baseline, and ensemble generation for all 50 states is computationally unnecessary.

\subsubsection{Reproducibility}

All algorithmic specifications, METIS parameters, and input data are publicly available:
\begin{itemize}
    \item \textbf{Source code}: github.com/[repository]/redistricting-algorithms
    \item \textbf{Census data}: 2020 Census tract boundaries and populations (Census Bureau)
    \item \textbf{METIS version}: 5.1.0 via \texttt{redist} Rust binary (C METIS FFI via \texttt{metis} crate v0.2)
    \item \textbf{Computational environment}: Rust 1.77+ (\texttt{redist} binary); post-hoc data analysis (CSV loading, plotting) uses Python 3.13 with \texttt{redist-analysis} crate (compactness, VRA, partisan metrics)
\end{itemize}

Independent researchers can replicate our algorithmic plans exactly by using identical METIS parameters and random seeds.
